% VI21-FULL-SOLUTION-REFINEMENT: P10
\subsection*{VI.21.P10: $k$-rational points of an affine $k$-scheme}
\begin{problem}\label{prob:vi21-p10}
Let $X=\Spec A$ for a $k$-algebra $A$.  Prove
\[
X(k)=\operatorname{Hom}_{\Spec k}(\Spec k,X)
\cong\operatorname{Hom}_{k\text{-alg}}(A,k).
\]
\end{problem}

\ifdefined\FullProblemDossiers
\paragraph{Why this problem matters.}
Relative rational points are affine ring maps that also respect the chosen base.

\paragraph{Useful ideas before starting.}
Translate point statements through local rings and residue fields; when the target is affine, use the VI/19 universal property and then verify the pointwise interpretation.

\paragraph{Guided hints.}
\begin{enumerate}
\item Identify the relevant canonical ring or stalk homomorphism.
\item Track maximal ideals or prime contractions to identify the image point.
\item Use locality to pass to residue fields when needed.
\item Check independence, uniqueness, or functoriality explicitly rather than only on points.
\end{enumerate}

\begin{solution}
The $k$-algebra structure on $A$ is a ring map $k\to A$, which contravariantly defines the
structure morphism $X=\Spec A\to\Spec k$.  A morphism $u:\Spec k\to X$ over $\Spec k$ corresponds
by affine anti-equivalence to a ring map $\alpha:A\to k$.  The triangle over $\Spec k$ commutes
exactly when the reversed ring triangle
\[
k\longrightarrow A\xrightarrow{\alpha}k
\]
is the identity on $k$.  That is precisely the condition that $\alpha$ be a $k$-algebra
homomorphism.  Conversely every $k$-algebra homomorphism gives such an over-$k$ morphism.  Hence the
bijection is natural.
\end{solution}

\paragraph{What happened mathematically.}
Relative rational points are affine ring maps that also respect the chosen base.

\paragraph{Extensions.}
Replace $k$ by an extension field $K$ and obtain $X(K)\cong\operatorname{Hom}_{k	ext{-alg}}(A,K)$.

\paragraph{Provenance.}
\textbf{Canonical VI/21 derivative/application; not an additional legacy migration.}
\else
\paragraph{Provenance.} \textbf{Canonical VI/21 derivative/application; not an additional legacy migration.}
\fi
